%% EXHALE: radiative cooling of a helium-dominated wind.
%% What cools a gas that has almost no hydrogen: the line channels available
%% to He, the thresholds implemented in Cool_coeff.f90, the numerical size of
%% each against Lyman-alpha, and the one channel the code does not carry.
%% Written for the LHS 1140b lower-atmosphere extension
%% (docs/lhs1140b_lower_atmosphere_plan_new.md, Phase C).
%% Author: Kwang-Il Seon
\documentclass[english,a4paper]{my_memo}
\synctex=1
\usepackage{booktabs}
\usepackage{amsmath}
\usepackage{graphicx}
\usepackage{babel}
\emergencystretch=3em

\title{Radiative cooling of a helium-dominated wind}
\author{Kwang-Il Seon}
\date{Last updated: \docmoddate}

\begin{document}
\maketitle

\begin{abstract}
A helium-dominated escaping atmosphere has almost no line cooling. Every
line channel available to helium has a high excitation threshold
($19.8$~eV for the lowest, against $10.2$~eV for Lyman-$\alpha$), so the
Boltzmann factor suppresses helium line cooling by five to six orders of
magnitude relative to Lyman-$\alpha$ at $10^4$~K. The energy budget of such
a wind therefore closes on continuum channels and adiabatic expansion rather
than on lines, and the gas has no effective thermostat. This memo records
the channels implemented in \texttt{Cool\_coeff.f90}, their measured
relative sizes, the one physically real channel the code does not carry
(the He~\textsc{i} $584$~\AA{} resonance line), and what
Cherubim et al. (2026) do and do not say about any of this.
\end{abstract}

\section{Why the question arises}

The metastable helium detection at LHS~1140b \cite{Cherubim2026}
(hereafter C26) is interpreted as a wind with
$\mathrm{H\!:\!He} \lesssim 10^{-3}$, that is, a gas in which hydrogen is a
trace species. Every EXHALE configuration validated to date is
hydrogen-dominated, and in all of them Lyman-$\alpha$ is the dominant line
coolant. Removing hydrogen therefore removes the coolant the code has always
relied on, and the question of what replaces it must be answered before a
helium-rich wind solution can be trusted. This is the energy-closure item of
Phase~C in \texttt{lhs1140b\_lower\_atmosphere\_plan\_new.md}.

\section{Line channels available to helium}

\subsection{What the code implements}

The collisional-excitation cooling terms are assembled in
\texttt{util\_ion\_eq.f90} and their coefficients defined in
\texttt{Cool\_coeff.f90}. Reading the exponential threshold of each
coefficient identifies the transition it represents:

\begin{table}[h]
\centering
\small
\begin{tabular}{@{}llll@{}}
\toprule
Coefficient & Threshold & Transition & Routine \\
\midrule
$7.5\!\times\!10^{-19}(1+\sqrt{T/10^5})^{-1}e^{-118348/T}$
  & $10.20$~eV & H\,\textsc{i} Ly$\alpha$ $1216$~\AA{}
  & \texttt{coex\_rate\_HI} \\
$1.1\!\times\!10^{-19}T^{0.082}e^{-2.3\times10^{5}/T}$
  & $19.82$~eV & He\,\textsc{i} $1^1S \rightarrow 2^3S$
  & \texttt{coex\_rate\_HeI} \\
$5.54\!\times\!10^{-17}T^{-0.397}(1+\sqrt{T/10^5})^{-1}e^{-473638/T}$
  & $40.81$~eV & He\,\textsc{ii} Ly$\alpha$ $304$~\AA{}
  & \texttt{coex\_rate\_HeII} \\
$1.16\!\times\!10^{-20}\sqrt{T}\,e^{-13179/T}$
  & $1.14$~eV & He\,\textsc{i} $2^3S \rightarrow 2^3P$, $10830$~\AA{}
  & \texttt{coex\_rate\_HeI23S\_10830} \\
\bottomrule
\end{tabular}
\caption{Collisional-excitation cooling coefficients (cgs, to be multiplied
by $n_e$ and the population of the lower level). The threshold energies are
the exponential arguments converted with $k_{\rm B}$; each matches a known
level to three digits, which is how the transitions above were identified
rather than assumed.}
\end{table}

Two further terms act on the metastable population when the triplet is
tracked: the collisional conversions $2^3S \rightarrow 2^1S$ ($0.80$~eV) and
$2^3S \rightarrow 2^1P$ ($1.40$~eV), through
\texttt{coex\_HeI\_23S\_21S} and \texttt{coex\_HeI\_23S\_21P}. The
ground-to-triplet excitation $q_{13}$ is deliberately excluded there because
the $19.82$~eV term in the table already includes it; the code comment
records this to prevent counting it twice.

The $10830$~\AA{} channel and the two conversions scale with the explicitly
computed $n(2^3S)$. That population is supplied by recombination and is
small, so these channels track a minority level and cannot act as a bulk
thermostat, however low their thresholds are.

\subsection{What the single exponential actually represents}

The $584$~\AA{} resonance line ($1^1S \rightarrow 2^1P$, $21.22$~eV) is the
optically allowed counterpart of Lyman-$\alpha$ and has no term of its own in
\texttt{Cool\_coeff.f90}. Reading the thresholds alone, one would conclude
that the code includes only the $19.82$~eV metastable channel and leaves out
the resonance line. Checking that conclusion against CHIANTI shows it is wrong,
in both directions, and the correction matters enough to state explicitly.

Summing every transition out of the He\,\textsc{i} ground level in CHIANTI~v11
(\texttt{cooling\_data/fit\_helium\_cooling.py}) gives the following. First,
the $584$~\AA{} line is \emph{not} the dominant channel at wind
temperatures: relative to the $2^3S$ channel it contributes $1\%$ at
$5000$~K, $6\%$ at $10^4$~K and $52\%$ at $5\times10^4$~K, because its
$1.4$~eV higher threshold costs more than its larger collision strength
returns until the gas is very hot. Second, and more usefully, the existing
coefficient is not a fit to the $2^3S$ line at all: it reproduces the
\emph{total} ground-state sum to within $5$--$19\%$ over
$5\times10^3$--$5\times10^4$~K, while departing from the $2^3S$ channel alone
by up to a factor $3.4$. Its $19.82$~eV exponent gives it the temperature
dependence of the lowest channel, and its prefactor is large enough to match
the sum of all of them.

The resonance line is therefore already included in the total this formula
was fitted to, so the practical question is how accurate that formula is, not
whether the line is present. Section~\ref{sec:chianti}
replaces both helium coefficients with CHIANTI fits and reports what that is
worth; the larger error turns out to be in He\,\textsc{ii}, not He\,\textsc{i}.

\section{How weak the helium channels are}

Evaluating the coefficients of Table~1 directly:

\begin{table}[h]
\centering
\begin{tabular}{lllll}
\toprule
$T$ [K] & H\,\textsc{i} Ly$\alpha$ & He\,\textsc{i} $19.8$~eV
        & He\,\textsc{ii} $304$~\AA{} & He\,\textsc{i}/H\,\textsc{i} \\
\midrule
5000 & $3.22\times10^{-29}$ & $2.33\times10^{-39}$ & $1.12\times10^{-59}$ & $7.2\times10^{-11}$ \\
8000 & $2.20\times10^{-25}$ & $7.51\times10^{-32}$ & $2.36\times10^{-44}$ & $3.4\times10^{-7}$ \\
10000 & $4.13\times10^{-24}$ & $2.40\times10^{-29}$ & $2.93\times10^{-39}$ & $5.8\times10^{-6}$ \\
15000 & $2.02\times10^{-22}$ & $5.30\times10^{-26}$ & $1.70\times10^{-32}$ & $2.6\times10^{-4}$ \\
20000 & $1.40\times10^{-21}$ & $2.51\times10^{-24}$ & $3.90\times10^{-29}$ & $1.8\times10^{-3}$ \\
30000 & $9.38\times10^{-21}$ & $1.20\times10^{-22}$ & $8.31\times10^{-26}$ & $1.3\times10^{-2}$ \\
50000 & $4.12\times10^{-20}$ & $2.69\times10^{-21}$ & $3.40\times10^{-23}$ & $6.5\times10^{-2}$ \\
\bottomrule
\end{tabular}
\caption{Collisional-excitation cooling coefficients in cgs, evaluated from
the implemented expressions. The last column is the ratio of the
He\,\textsc{i} channel to Lyman-$\alpha$ at equal lower-level density.}
\end{table}

The comparison is per atom of the relevant species, so in a wind where
helium outnumbers hydrogen by $10^3$ the ratio in the last column is raised
by that factor; even then the helium channel only reaches parity with what a
trace hydrogen population would supply near $10^4$~K, and falls behind at
lower temperature. The He\,\textsc{ii} line is irrelevant below
$\sim 3\times10^4$~K: at $40.81$~eV it is essentially extinct in a
photoionized wind at a few thousand kelvin.

\section{What closes the energy budget instead}

With lines suppressed, the terms that survive are the continuum ones already
in \texttt{util\_ion\_eq.f90} --- recombination cooling of
H\,\textsc{ii}/He\,\textsc{ii}/He\,\textsc{iii}, collisional-ionization
cooling, and free-free emission --- together with the $P\,dV$ work of the
expanding wind, which is not a radiative term at all but is what actually
carries the energy away in a fast outflow.

The first converged case of the Phase~C sweep already shows the transition.
At $\mathrm{He/H} = 1$ (the tutorial planet with only the helium ratio
changed), the channel-resolved output gives heating led by He\,\textsc{i}
($96\%$ of the total at the base) and cooling led by recombination
($75\%$ at the base, $91\%$ at the top); Lyman-$\alpha$ leads only in the
intermediate layer at $1.5\,R_{\rm p}$ ($54\%$, $T = 1.6\times10^4$~K),
where hydrogen is still half of the nuclei. The remaining cases of the sweep
($\mathrm{He/H} = 10$, $100$, $1000$) test whether that last
Lyman-$\alpha$ peak disappears and the budget closes on the continuum
channels alone, which is the acceptance criterion for this phase.

\section{CHIANTI-based replacement coefficients}
\label{sec:chianti}

Both helium coefficients were refitted to CHIANTI~v11 with the same method
already used for the metal coolants
(\texttt{cooling\_data/chianti\_cooling.py} reads the raw
\texttt{.elvlc}/\texttt{.wgfa}/\texttt{.scups} files and descales the
effective collision strengths following Burgess \& Tully 1992; no ChiantiPy).
The fitted quantity is the sum over all transitions out of the ground level,
which is what the existing coefficients represent, so the rule of Section~2.1 is unchanged: the
ground-to-triplet channel stays inside this term, and is still left out of the
triplet system so that it is not counted twice.

The functional form is the one \texttt{fit\_cooling\_formulas.py} uses, exact
for a single line and a good effective form for a sum whose shape is set by
its lowest threshold,
%
\begin{equation}
  \Lambda(T) = \frac{C}{\sqrt{T}}
               \left[a + b\ln\!\left(1 + \frac{T}{T_0}\right)\right]
               e^{-T_{\rm ex}/T},
  \qquad C = \frac{8.629\times10^{-6}\,\Delta E}{g_l},
  \qquad T_{\rm ex} = \frac{\Delta E}{k_{\rm B}},
\end{equation}
%
with $\Delta E$ left free so the fit recovers the effective threshold of the
sum rather than assuming the lowest level. It recovers $19.88$~eV for
He\,\textsc{i} and $40.76$~eV for He\,\textsc{ii}, both within $0.3\%$ of the
lowest level of each ion, which is the expected behavior for a
threshold-dominated sum.

\begin{table}[h]
\centering
\small
\begin{tabular}{@{}lllllll@{}}
\toprule
Ion & $g_l$ & $C$ & $a$ & $b$ & $T_0$ [K] & $T_{\rm ex}$ [K] \\
\midrule
He\,\textsc{i}  & 1 & $2.749039\times10^{-16}$ & $0.0613017$ & $0.4191808$
                & $9.821387\times10^{4}$ & $230747.5$ \\
He\,\textsc{ii} & 2 & $2.817309\times10^{-16}$ & $0.4371421$ & $0.6051062$
                & $1.905585\times10^{5}$ & $472956.0$ \\
\bottomrule
\end{tabular}
\caption{Fitted coefficients, cgs. Fit ranges $3\times10^3$--$10^5$~K
(He\,\textsc{i}) and $10^4$--$2\times10^5$~K (He\,\textsc{ii}).}
\end{table}

Table~4 gives the accuracy against the CHIANTI sum, together with the error
of the coefficients these would replace.

\begin{table}[h]
\centering
\small
\begin{tabular}{@{}lll@{}}
\toprule
 & this fit & current \texttt{Cool\_coeff.f90} \\
\midrule
He\,\textsc{i}, $5\times10^3$--$5\times10^4$~K
  & $0.4\%$ & $-19.4\%$ (ratio $0.81$--$0.95$) \\
He\,\textsc{ii}, $2\times10^4$--$10^5$~K
  & $0.2\%$ & $-53.6\%$ (ratio $0.46$--$0.73$) \\
\bottomrule
\end{tabular}
\caption{Maximum deviation from the CHIANTI ground-state sum over the range
where each coolant can matter.}
\end{table}

The He\,\textsc{i} coefficient is therefore adequate at the tens-of-percent
level and biased low; the He\,\textsc{ii} coefficient is low by a factor of
two and degrades with temperature, so it is the one worth replacing on
physical grounds. Neither error is important in a hydrogen-rich run, where
Lyman-$\alpha$ exceeds both by orders of magnitude, which is presumably why
they survived; both act directly on the thermal structure of a
helium-dominated wind. Fitting scripts, coefficient file
(\texttt{helium\_cooling\_fits.txt}) and comparison figure
(\texttt{helium\_cooling\_fits.pdf}) are in \texttt{cooling\_data/}. The
coefficients are recorded here but not yet installed in the Fortran; doing so
changes converged solutions and therefore needs its own golden refresh.

A caveat that applies to both, and to the coefficients they replace: this is
the optically thin, low-density (coronal) limit, in which every collisional
excitation is followed by a radiative decay that escapes. The $584$~\AA{}
line is a resonance line of the dominant species in a helium-rich wind, so
it can be optically thick, and the $2^3S$ channel feeds a metastable level
that EXHALE tracks explicitly when the triplet is on. Both effects reduce the
true cooling below the coronal value; the fits above bound it from that side
rather than resolving it.

\section{Consequence: a wind without a thermostat}

Line cooling is what keeps a hydrogen-rich thermosphere near
$10^4$~K: Lyman-$\alpha$ switches on sharply once the gas reaches the
excitation threshold and absorbs further heating at nearly constant
temperature. A helium-dominated gas has no such switch below several times
$10^4$~K. It should therefore run hotter for the same heating rate, and its
temperature should be set by the balance between photoionization heating,
recombination, and expansion rather than pinned by a line.

This is consistent with what the C26 retrieval requires
($T_{\rm wind} \simeq 5160$~K, with their grid explored to $6400$~K) and
with our own p-winds oracle, in which temperature is by far the strongest
lever on the line depth while the mass-loss rate is weak and saturated
(\texttt{LHS1140b/pwinds\_oracle/README.md}). It is a consistency argument,
not yet a demonstration: only a solution with a computed energy equation can
show that such a temperature is actually maintained.

\section{What C26 says about cooling}

Nothing. The paper contains no radiative-cooling discussion, no energy
equation, and no heating efficiency. This is structural rather than an
oversight: p-winds models the outflow as an \emph{isothermal} Parker wind,
so the outflow temperature is an input. The Supplement states it plainly:
the code forward models the spectrum ``for given values of the atmospheric
escape rate, outflow temperature, H:He atomic number ratio, line-of-sight
wind speed, top-of-atmosphere XUV spectrum, planetary mass and radius'', and
$T_{\rm wind}$ is then estimated by MCMC alongside the other three
parameters. The stellar spectrum enters only through the ionization rates
and the resulting number densities, never as a heating rate.

A full-text search of the article and its Supplement finds the word
``cool'' in three places, none of them about cooling physics: ``cooler,
rocky exoplanets'' in the introduction, ``the cool equilibrium temperature''
in the cold-trap discussion, and a reference title. ``Isothermal'' appears
twice: once for the wind and once for the APEC model of the stellar corona,
which concerns the X-ray fit and not the planet. No heating efficiency
$\eta$ is quoted anywhere.

The two places where energetics do appear are global budget arguments, not
statements about the interior of the wind:

\begin{enumerate}
\item an energy-limited estimate of $(6.2\!-\!29)\times10^{7}$~g\,s$^{-1}$,
      obtained by integrating the scaled SEDs over $10\!-\!1300$~\AA{},
      which brackets their retrieved
      $\dot{M} = 2.03^{+0.58}_{-0.67}\times10^{8}$~g\,s$^{-1}$;
\item the exclusion of core-powered mass loss on the grounds that it
      predicts a much smaller rate.
\end{enumerate}

A related check they do report is a composition test rather than an energy
one: fixing the mean molecular weight at $4$~amu, the theoretical maximum
for a pure neutral helium wind, leaves the transmission spectra unchanged.

\section{Implication for the extension}

The retrieved $T_{\rm wind}$ is an observational constraint on the line
profile and the metastable population, not a temperature shown to be
thermodynamically sustainable. Supplying that missing half is precisely the
value of running EXHALE on this target: EXHALE computes the temperature from
a heating--cooling balance, so the question ``does a helium-dominated wind
settle near $5000\!-\!6400$~K under this irradiation?'' becomes answerable
rather than assumed. The weakness of helium line cooling documented above
makes a hot, poorly regulated solution the physically expected outcome, and
makes the accuracy of the few surviving channels --- including the missing
$584$~\AA{} line --- more consequential than it would be in any
hydrogen-rich run.

\begin{thebibliography}{9}
\bibitem[Cherubim et al.\ (2026)]{Cherubim2026}
C.~Cherubim et al.,
\emph{Helium escaping from the atmosphere of a nearby rocky exoplanet
orbiting in a habitable zone},
Science, 10.1126/science.aea9708 (2026);
first release 2026 July 16, with Supplementary Materials.
\end{thebibliography}

\end{document}
